% IN5340/IN9340 Project I Presentation
% Build: pdflatex presentation_p1.tex

\documentclass[aspectratio=169,10pt]{beamer}

\usetheme{Madrid}
\usecolortheme{default}

\usepackage{graphicx}
\usepackage{amsmath,amssymb}
\usepackage{booktabs}
\usepackage{mathtools}
\usepackage[T1]{fontenc}
\usepackage[utf8]{inputenc}

\graphicspath{{figs/}}

% Figure helpers (keeps plots within the slide)
\newcommand{\fullfig}[2][]{%
	\centering
	\includegraphics[width=\linewidth,height=0.78\textheight,keepaspectratio,#1]{#2}%
}
\newcommand{\colfig}[2][]{%
	\includegraphics[width=\linewidth,height=0.58\textheight,keepaspectratio,#1]{#2}%
}

\title[IN5340 Project I]{IN5340/IN9340 Project I: Random Variables and Random Processes}
\author{Theodor W\aa lberg}
\institute{University of Oslo}
\date{February 2026}

\begin{document}

\begin{frame}
	\titlepage
	\vspace{-1.0em}
	\centering
	\footnotesize
	Code: \href{https://github.com/Twaal/Statistical_Signal_Processing_IN5340/tree/master/Project_1}{\nolinkurl{github.com/Twaal/Statistical_Signal_Processing_IN5340/tree/master/Project_1}}
\end{frame}

\begin{frame}{Agenda}
	\tableofcontents
\end{frame}

\section{1. Random Variables}

\subsection{1.1 Uniform random variable \texorpdfstring{$X \sim \mathcal{U}(0,1)$}{X ~ U(0,1)}}

\begin{frame}{Uniform RV: Task}
\begin{itemize}
	\item Generate a long sequence $x[n]$ of i.i.d. samples from $\mathcal{U}(0,1)$.
	\item Estimate the PDF using a manually-normalized histogram.
	\item Estimate mean $m_x = \mathbb{E}[X]$ and variance $\sigma_x^2 = \mathbb{E}[(X-m_x)^2]$.
	\item Overlay the theoretical PDF and compare theory vs estimates.
\end{itemize}
\vspace{0.5em}
\small
Simulation in notebook: $N=5000$ samples (seeded).
\end{frame}

\begin{frame}{Uniform RV: Theoretical mean and variance derivation}
Assume $f_X(x)=1$ for $x\in[0,1]$ and $0$ otherwise.
\vspace{0.5em}
\begin{align*}
\mathbb{E}[X]
&= \int_{-\infty}^{\infty} x\, f_X(x)\,dx
= \int_{0}^{1} x\,dx
= \left.\frac{x^2}{2}\right|_{0}^{1}
= \frac{1}{2}.\\
\mathbb{E}[X^2]
&= \int_{0}^{1} x^2\,dx
= \left.\frac{x^3}{3}\right|_{0}^{1}
= \frac{1}{3}.\\
\mathrm{Var}(X)
&= \mathbb{E}[X^2]-\mathbb{E}[X]^2
= \frac{1}{3}-\left(\frac{1}{2}\right)^2
= \frac{1}{12}\approx 0.0833.
\end{align*}
\end{frame}

\begin{frame}{Uniform RV: Sequence and basic statistics}
\begin{columns}[T,onlytextwidth]
\column{0.62\textwidth}
\centering
\colfig{uniform_rand_seq.png}
\column{0.38\textwidth}
\small
\textbf{Estimated (one run):}
\begin{itemize}
	\item $\hat m_x \approx 0.4968$
	\item $\hat\sigma_x^2 \approx 0.08387$
\end{itemize}
\textbf{Theory:}
\begin{itemize}
	\item $m_x = 1/2$
	\item $\sigma_x^2 = 1/12 \approx 0.08333$
\end{itemize}
\end{columns}
\end{frame}

\begin{frame}{Uniform RV: PDF estimate vs theory}
\fullfig{uniform_histogram_vs_theory.png}
\vspace{0.5em}
\small
Histogram is normalized to integrate to 1; the theoretical PDF is constant on $[0,1]$.
\end{frame}

\subsection{1.2 Gaussian random variable \texorpdfstring{$X \sim \mathcal{N}(\mu,\sigma^2)$}{X ~ N(mu,sigma^2)}}

\begin{frame}{Gaussian RV: Task}
\begin{itemize}
	\item Generate i.i.d. samples from a Gaussian RV (Matlab: \texttt{randn}, Python: \texttt{np.random.normal}).
	\item Estimate PDF (histogram) + mean and variance.
	\item Overlay the theoretical Gaussian PDF.
\end{itemize}
\vspace{0.5em}
\small
In the notebook: $\mu=0$, $\sigma=1$, $N=5000$.
\end{frame}

\begin{frame}{Gaussian RV: Theoretical mean and variance}
For $X\sim\mathcal{N}(\mu,\sigma^2)$ with density
\[
f_X(x)=\frac{1}{\sqrt{2\pi\sigma^2}}\exp\left(-\frac{(x-\mu)^2}{2\sigma^2}\right),
\]
the moments are
\[
\mathbb{E}[X]=\mu,\qquad \mathrm{Var}(X)=\sigma^2.
\]
\vspace{0.5em}
\small
Here: $\mu=0$, $\sigma^2=1$.
\end{frame}

\begin{frame}{Gaussian RV: Sequence}
\begin{columns}[T,onlytextwidth]
\column{0.62\textwidth}
\centering
\colfig{normal_random_seq.png}
\column{0.38\textwidth}
\small
\textbf{Estimated (one run):}
\begin{itemize}
	\item $\hat\mu \approx -2.2\cdot10^{-4}$
	\item $\widehat{\sigma^2} \approx 1.0013$
\end{itemize}
\textbf{Theory:}
\begin{itemize}
	\item $\mu=0$
	\item $\sigma^2=1$
\end{itemize}
\end{columns}
\end{frame}

\begin{frame}{Gaussian RV: PDF estimate vs theory}
\fullfig{normal_histogram_vs_theory.png}
\vspace{0.5em}
\small
Histogram-based density estimate with the theoretical Gaussian curve overlay.
\end{frame}

\begin{frame}{Gaussian RV: Experimental vs theoretical PDF}
\fullfig{normal_curve_vs_theory.png}
\end{frame}

\subsection{1.3 Central Limit Theorem}

\begin{frame}{CLT: Setup}
\begin{itemize}
	\item Generate 12 independent uniform RVs and average:
	\[
	Y = \frac{1}{12}\sum_{k=1}^{12} X_k,\qquad X_k \overset{iid}{\sim} \mathcal{U}(0,1).
	\]
	\item Estimate the PDF of $Y$ and compare to a Gaussian approximation.
	\item Derive theoretical mean and variance for $Y$.
\end{itemize}
\end{frame}

\begin{frame}{CLT: Theoretical mean and variance derivation}
Using linearity of expectation and independence:
\begin{align*}
\mathbb{E}[Y]
&= \mathbb{E}\left[\frac{1}{12}\sum_{k=1}^{12}X_k\right]
= \frac{1}{12}\sum_{k=1}^{12}\mathbb{E}[X_k]
= \mathbb{E}[X]
= \frac{1}{2}.\\
\mathrm{Var}(Y)
&= \mathrm{Var}\left(\frac{1}{12}\sum_{k=1}^{12}X_k\right)
= \frac{1}{12^2}\sum_{k=1}^{12}\mathrm{Var}(X_k)\quad (\text{independent})\\
&= \frac{12}{144}\cdot\frac{1}{12}
= \frac{1}{144}\approx 0.00694.
\end{align*}
\vspace{0.25em}
\small
In the notebook, $Y$ is centered as $Y_\text{norm}=Y-\mathbb{E}[Y]$, so $\mathbb{E}[Y_\text{norm}]=0$ and $\mathrm{Var}(Y_\text{norm})=\mathrm{Var}(Y)$.
\end{frame}

\begin{frame}{CLT: Histogram of sample means}
\begin{columns}[T,onlytextwidth]
\column{0.62\textwidth}
\centering
\colfig{clt_histogram_of_sample_means.png}
\column{0.38\textwidth}
\small
\textbf{Estimated (one run):}
\begin{itemize}
	\item $\hat m \approx 0$ (after centering)
	\item $\hat\sigma^2 \approx 0.00693$
\end{itemize}
\textbf{Theory:}
\begin{itemize}
	\item $m = 0$
	\item $\sigma^2 = 1/144 \approx 0.00694$
\end{itemize}
\end{columns}
\end{frame}

\section{2. Stationarity and Ergodicity}

\begin{frame}{Random processes in the assignment}
We analyze three random processes (generated from pseudo-random uniform noise):
\vspace{0.25em}
\small
\begin{align*}
\text{rp1: }\; v[n] &= \big(\mathrm{rand}-0.5\big)\,\underbrace{b\sin(\pi n/N)}_{\text{time-varying scale}} + \underbrace{a n}_{\text{time-varying mean}},\quad \\
\text{rp2: }\; v[n] &= \big(\mathrm{rand}-0.5\big)\,m + a,\\
\text{rp3: }\; v[n] &= \big(\mathrm{rand}-0.5\big)\,M_r + A_r,\quad A_r,M_r \text{ fixed per realization.}
\end{align*}
\vspace{0.25em}
\normalsize
Goal: decide stationarity/ergodicity by inspection and by ensemble/time averages.
\end{frame}

\begin{frame}{2.1 Inspection: realizations (M=4, N=100)}
\fullfig{random_processes.png}
\end{frame}

\begin{frame}{2.1 Stationary vs ergodic (qualitative conclusion)}
\begin{itemize}
	\item \textbf{Process 1}: non-stationary (mean and scale vary with time) $\Rightarrow$ non-ergodic.
	\item \textbf{Process 2}: stationary (constant mean/variance) and appears ergodic in mean/variance.
	\item \textbf{Process 3}: stationary in distribution, but realization-dependent offset/scale $\Rightarrow$ non-ergodic.
\end{itemize}
\end{frame}

\begin{frame}{2.2 Ensemble averages (M=80, N=200)}
\fullfig{ensemble_stats.png}
\vspace{0.25em}
\small
Ensemble mean/standard deviation vs time: rp1 varies with time; rp2 roughly constant; rp3 roughly constant but hides realization-dependent parameters.
\end{frame}

\begin{frame}{2.3 Time averages + images (M=4, N=1000)}
\fullfig{ergodicity.png}
\vspace{0.25em}
\small
Imagesc-style view: time on $x$, realization on $y$. Time-variation across columns suggests non-stationarity; large realization-to-realization offsets suggest non-ergodicity.
\end{frame}

\begin{frame}{2.3 Time averages: example output (one run)}
\small
\begin{itemize}
	\item rp1 time-averages (per realization) are similar but the process is non-stationary, so not ergodic.
	\item rp2 time-averages cluster around the ensemble mean (consistent with ergodicity).
	\item rp3 time-averages differ significantly across realizations (non-ergodic in mean).
\end{itemize}
\vspace{0.5em}
\footnotesize
Example (from notebook stdout):
\[
\bar x_{\text{time}}^{(rp2)} \approx [0.475,\;0.453,\;0.526,\;0.515],\qquad
\bar x_{\text{time}}^{(rp3)} \approx [0.154,\;0.909,\;0.387,\;0.047].
\]
\end{frame}

\section*{Summary}

\begin{frame}{Summary}
\begin{itemize}
	\item Uniform and Gaussian RV experiments: estimated mean/variance and PDF match theory with small finite-sample deviations.
	\item CLT: the average of 12 uniforms is well-approximated by a Gaussian; theory gives $\sigma^2=1/144$.
	\item Random processes: rp1 non-stationary/non-ergodic; rp2 stationary/ergodic; rp3 stationary but non-ergodic.
\end{itemize}
\vspace{0.5em}
\small
All code and figures reproduced from the accompanying notebook.
\end{frame}

\end{document}
